\documentclass[11pt]{article}
\usepackage[margin=1in]{geometry}
\usepackage[T1]{fontenc}
\usepackage{lmodern,amsmath,amssymb,amsthm,mathtools,booktabs,array,xcolor}
\usepackage{hyperref}
\hypersetup{colorlinks=true,linkcolor=blue!45!black,citecolor=blue!45!black,urlcolor=blue!45!black,
pdfauthor={Joshua Mark},pdftitle={Beyond asymptotic Hadamard existence: arithmetic and geometric research directions}}
\setlength{\emergencystretch}{3em}
\newtheorem{proposition}{Proposition}[section]
\newtheorem{lemma}[proposition]{Lemma}
\newtheorem{question}[proposition]{Research question}
\newtheorem{remark}[proposition]{Remark}
\newcommand{\Z}{\mathbb Z}
\newcommand{\R}{\mathbb R}
\newcommand{\N}{\mathbb N}
\newcommand{\one}{\mathbf 1}
\title{Beyond asymptotic Hadamard existence:\\arithmetic and geometric research directions}
\author{Joshua Mark\thanks{ORCID: \href{https://orcid.org/0009-0003-5905-2382}{0009-0003-5905-2382}. AI tools assisted with literature research, mathematical exploration, experiments and exposition.}}
\date{20 November 2026 --- working research notebook}
\begin{document}
\maketitle
\begin{abstract}
We investigate routes to exact Hadamard existence independent of sparse
multibase representations. The most concrete arithmetic target is a
congruence-constrained Hermitian norm problem in integral cyclic group
rings: at its minimum energy, parity forces every coefficient to be a
sign. A second route uses integer quadratic moment semigroups and their
holes; a third uses low-rank integral Gram completion and row exchanges.
We prove elementary reductions and obstructions, construct exact finite
certificates, and identify several invalid strengthening attempts.
No all-order existence theorem, new universal construction, or proof of
the Hadamard conjecture is obtained. Open targets are explicitly separated
from established facts and finite experiments. No Lean work is undertaken.
\end{abstract}

\section{Objective and source status}
The target is
\[
\forall n\ge1,\quad 4\mid n\ \Longrightarrow\
\exists H\in\{\pm1\}^{n\times n}:\ HH^{\mathsf T}=nI_n.
\]
This notebook does not assume a particular construction ansatz. It is
separate from the existing asymptotic-existence project. The October
2026 survey~\cite{Survey} still presents the universal statement as open;
its finite-order table is not a proof of the universal assertion. Public
source searches were checked again on 20 November 2026. This is a
targeted investigation, not an exhaustive priority audit.

We allow stronger theorems in other subjects as intermediate goals.
However, an equivalent reformulation is not progress on existence, and
a stronger statement may be false even if the Hadamard conjecture is
true. We therefore ask, for each candidate: what implication is proved,
what genuinely new mathematical step remains, and what counterexamples
must that step survive?

\section{Arithmetic norm representations at minimum energy}
\subsection{A coefficient bound that need not be assumed}
Let $m$ be odd and set
\[
 R_m=\Z[z]/(z^m-1),\quad z^*=z^{-1},\quad
 J_m=1+z+\cdots+z^{m-1}.
\]
Represent every ring element uniquely with degree below $m$. For
$A=\sum_j a_jz^j$, the coefficient of $1$ in $AA^*$ is $\sum_j a_j^2$.

\begin{proposition}[Parity and minimum energy]\label{prop:energy}
Suppose $A_1,\ldots,A_4\in R_m$ satisfy
\begin{equation}\label{eq:norm}
 A_i\equiv J_m\pmod{2R_m},\qquad
 \sum_{i=1}^4 A_iA_i^*=4t\cdot1.
\end{equation}
Then $t\ge m$ and $t$ is odd. If $t=m$, every coefficient of every
$A_i$ is $+1$ or $-1$.
\end{proposition}
\begin{proof}
All $4m$ coefficients are odd. Comparing constant coefficients gives
$\sum_{i,j}a_{ij}^2=4t\ge4m$. Equality forces every square to be one.
The augmentation map $\varepsilon(A)=A(1)$ gives
$\sum_i\varepsilon(A_i)^2=4t$. Since $m$ is odd, all four augmentations
are odd, so the left side is $4$ modulo $8$. Hence $t$ is odd.
\end{proof}

Thus one can seek an integral norm theorem with a congruence condition,
rather than impose exponentially many separate sign conditions.
The difficult part has not disappeared: exact integral solutions at
the minimum energy are still needed.

\begin{question}[A broader arithmetic lifting problem]\label{q:norm}
For every odd $m\ge1$ and every odd $t\ge m$, does
\eqref{eq:norm} have a solution?
\end{question}
This question is deliberately stronger than the boundary case needed
below. We have not proved that its elementary necessary conditions are
sufficient, or even excluded further arithmetic obstructions. It is a
testable question about a family of integral quadratic equations, not an
asserted theorem. A useful first advance would be a sufficient condition
for $t>m$ with explicit dependence on $m$, followed by an analysis of the
obstructions as $t$ decreases to $m$. A theorem that works only far above
$m$ would not by itself settle Hadamard existence.

\subsection{The precise implication to Hadamard matrices}
For a sign sequence $a=(a_0,\ldots,a_{m-1})$, write
\[
 P_a(s)=\sum_{j=0}^{m-1}a_ja_{j+s\bmod m}.
\]
The ring identity at $t=m$ says
\begin{equation}\label{eq:scalar}
 \sum_{i=1}^4P_{a_i}(s)=0\qquad(1\le s<m).
\end{equation}
Let $A,B,C,D$ now denote the corresponding circulant sign matrices, and
let $R$ be the permutation matrix $R_{ij}=1$ when $i+j=0\pmod m$.
The Goethals--Seidel construction~\cite[Theorem 3]{Survey} is
\begin{equation}\label{eq:gs}
 \begin{pmatrix}
 A&BR&CR&DR\\
 -BR&A&D^{\mathsf T}R&-C^{\mathsf T}R\\
 -CR&-D^{\mathsf T}R&A&B^{\mathsf T}R\\
 -DR&C^{\mathsf T}R&-B^{\mathsf T}R&A
 \end{pmatrix}.
\end{equation}
Its Gram matrix is $4mI$ under~\eqref{eq:scalar}. This uses circulant
commutativity and $RXR=X^{\mathsf T}$, and does not require the four
circulants to be symmetric or pairwise amicable.

\begin{proposition}[Conditional route to the full conjecture]
If \eqref{eq:norm} is solvable at $t=m$ for every odd $m$, then the
Hadamard conjecture holds. In particular an affirmative answer to
Question~\ref{q:norm} would imply it.
\end{proposition}
\begin{proof}
Proposition~\ref{prop:energy} and \eqref{eq:gs} give order $4m$ for each
odd $m$. Every positive multiple of four is $4\cdot2^a m$ for some
odd $m$, and $a$ Sylvester doublings give that order.
\end{proof}

The boundary assertion is an all-odd-length four-circulant existence
problem, a stronger construction claim than unrestricted Hadamard
existence. Writing it in ring language alone does not prove it.
The potentially useful import is integral representation theory with
congruences and controlled energy.
Hermitian methods already occur in generalized Hadamard existence and
completion~\cite{Brock}; the connection itself is not claimed as new.

\subsection{Why several attractive shortcuts fail}
At the trivial character $z=1$, the norm identity only requires a sum
of four odd squares equal to $4t$. Other characters must satisfy the
same identity simultaneously, with compatible integral coefficients.
For example at $m=5$ the four sequences
\[
 a=b=(1,1,1,1,-1),\qquad c=d=(1,1,1,-1,-1)
\]
have row sums $(3,3,1,1)$, whose squares sum to $20$. Their summed
autocorrelations at shifts $1,2$ are $4,-4$, so they fail
\eqref{eq:scalar}. Lagrange's scalar four-square theorem is not the
missing simultaneous representation theorem.

Likewise, solving the norm equation over $\R[C_m]$, or character by
character, does not provide the parity and integrality constraints.
Classical positive spectral factorization permits coefficients outside
the required arithmetic class. Perfect-reconstruction filter banks are
a useful language for polynomial norm identities~\cite{Filters}, but
their usual factorization freedom does not supply a binary solution.

There is a more decisive warning about quaternions. Put
$q_j=(a_j+b_j\mathbf i+c_j\mathbf j+d_j\mathbf k)/2$.
Equation~\eqref{eq:scalar} forces only the \emph{real parts} of
$\sum_jq_j\overline{q_{j+s}}$ to vanish. Setting the full quaternion
correlation to zero adds three constraints. The half-unit alphabet is
$q\mathcal Q_8\cup q^2\mathcal Q_8$, where
$q=(1+\mathbf i+\mathbf j+\mathbf k)/2$; multiplication by $q^{-1}$
sends it to $\mathcal Q_+=\mathcal Q_8\cup q\mathcal Q_8$ and preserves
perfection. Bennett--Bright--Colinot--Nayak~\cite{Quaternion} rule out
perfect $\mathcal Q_+$ sequences at lengths $35,47,53,59$, using their
equivalence theorem and earlier nonexistence results. Therefore that
stronger quaternionic statement is false, and cannot prove all orders.
Our Hermitian sum of four norms does not impose those extra conditions.

\subsection{Exact finite tests}
The supplied code finds four binary sequences satisfying
\eqref{eq:scalar} for $m=3,5,7,9$, and verifies the full Gram matrix of
each resulting order-$4m$ matrix. It also finds solutions to
\eqref{eq:norm} for $t=m+2$ at the same four $m$. In the latter case,
one coefficient has magnitude three and all others have magnitude one,
as the constant coefficient requires. These eight witnesses test only
small cases, not the general lifting assertion. The length-five binary
witness has quaternion correlation numerator $(0,0,-4,-4)$ at shift
one, with denominator four: it exhibits exactly why scalar cancellation
must not be strengthened inadvertently.

\section{Algebraic statistics: integer moment semigroups}
\subsection{Pairwise independence at the smallest possible support}
For $z\in\{0,1\}^d$, define its quadratic moment column
\[
 \phi(z)=(1,(z_i)_i,(z_iz_j)_{i<j}),\qquad
 b_{N,d}=(N,(N/2)_i,(N/4)_{i<j}).
\]
For $4\mid N$, a nonnegative integer solution of
\begin{equation}\label{eq:mom}
 \sum_z c_z\phi(z)=b_{N,d}
\end{equation}
is a multiset of $N$ binary observations whose coordinates are unbiased
and pairwise independent. Replacing bits by signs and adjoining an
all-one column gives $d+1$ orthogonal columns. Hence $N\ge d+1$.
For $d=N-1$ the result is precisely a Hadamard matrix of order $N$.
This tight-support boundary distinguishes the problem from approximate
covariance matching or a polynomial-size orthogonal array.

The general rare-object framework of Kuperberg--Lovett--Peled
\cite{KLP} constructs regular combinatorial objects with polynomial
overhead. Their discussion explicitly distinguishes minimal objects
and notes the difficulty of Gaussian counting for Hadamard designs.
The more recent partial-Hadamard counting theorem of Davis~\cite{Davis}
works when the number of columns is at least a constant times the cube
of the number of rows, not at square aspect ratio. Neither result
supplies a tight-support existence theorem.

\subsection{Signed feasibility is easy, positivity is not}
Let $e_U$ denote the unit vector at the bit pattern with support $U$.
\begin{proposition}[An explicit signed design]\label{prop:signed}
For every $d\ge2$ and $4\mid N$, the following integer coefficients
solve \eqref{eq:mom}:
\[
 c_\varnothing=\frac{N(d^2-5d+8)}8,\quad
 c_{\{i\}}=\frac{N(3-d)}4,\quad
 c_{\{i,j\}}=\frac N4,\quad
 c_U=0\quad(|U|\ge3).
\]
\end{proposition}
\begin{proof}
For any pair $i,j$ the only contributing pattern is $\{i,j\}$.
The first moment for $i$ is $N(3-d)/4+(d-1)N/4=N/2$.
Adding all coefficients gives $N$. The first coefficient is integral
because $d^2-5d+8$ is even.
\end{proof}

\begin{proposition}[All integer moment-preserving moves]\label{prop:moves}
For each $T\subseteq[d]$ with $|T|\ge3$, set
\[
 M_T=\sum_{U\subseteq T}(-1)^{|T|-|U|}e_U.
\]
These moves form a $\Z$-basis of the integer kernel of the quadratic
moment map. Consequently every integer solution of \eqref{eq:mom}
is the signed solution above plus $\sum_{|T|\ge3}a_TM_T$ for integers
$a_T$.
\end{proposition}
\begin{proof}
Every moment monomial has degree at most two, so its alternating sum
over a cube with at least three directions vanishes. Thus $M_T$ is in
the kernel. Its coefficient at $T$ is one and its other nonzero
coefficients are on strict subsets of $T$. Eliminate the coefficients
of any kernel vector in decreasing support size. A remaining vector
supported on sizes zero, one and two vanishes: pair moments first force
the size-two coefficients to zero, then first moments force singletons
to zero, and the total moment forces the empty coefficient to zero.
This also proves independence and integrality of the reduction.
\end{proof}

This gives a concrete positivity problem for integer cube moves, not
merely an analogy to algebraic geometry. In affine-semigroup language,
the target is in the nonnegative rational cone and in the generated
integer lattice. Such a point need not belong to the nonnegative
integer semigroup: it may be a \emph{hole}. The theory of cut polytopes
studies exactly this distinction~\cite{Ohsugi}.

For $d=N=4$, the signed solution has empty coefficient $2$, four
singleton coefficients $-1$, and six pair coefficients $1$. The
uniform fractional solution has coefficient $1/4$ at all sixteen
patterns. Yet no nonnegative integer solution exists: it would give
five orthogonal nonzero columns in $\R^4$. The script independently
exhausts all $\binom{19}{4}=3876$ four-row multisets. It also verifies
an eight-row positive solution. This is a genuine hole, not an example
at an admissible Hadamard boundary; the rank obstruction is intentional.
It disproves any unrestricted ``fractional plus lattice implies
integral'' bridge.

The worthwhile broader problem is to understand which holes survive
after spectral rank obstructions and symmetries are accounted for.
At $d=N-1$, asserting absence of the particular central holes is
equivalent to Hadamard existence, and we do not count it as a new theorem.
A general Markov-basis connectivity theorem also does not settle this:
connectivity of an already nonempty nonnegative fiber cannot establish
that the fiber is nonempty.

\section{Lattices, tight frames and completion}
Let $n$ be a positive multiple of four, $1\le r<n$, and let
$P\in\{\pm1\}^{(n-r)\times n}$ satisfy
$PP^{\mathsf T}=nI_{n-r}$. Define the residual Gram matrix
\[
 C=nI_n-P^{\mathsf T}P.
\]
\begin{proposition}[Exact residual criterion]\label{prop:residual}
The matrix $C$ is positive semidefinite, has rank $r$, satisfies
$C^2=nC$, and has diagonal entries $r$. Moreover $P$ completes to a
Hadamard matrix if and only if
\begin{equation}\label{eq:factor}
 C=B^{\mathsf T}B\quad\hbox{for some }B\in\{\pm1\}^{r\times n}.
\end{equation}
For even $n$, all off-diagonal entries of $C$ have parity $r$.
\end{proposition}
\begin{proof}
$P^{\mathsf T}P/n$ is an orthogonal projection of rank $n-r$; its
complement gives the assertions about $C$. Its diagonal is
$n-(n-r)=r$. Off-diagonal entries are negatives of sums of $n-r$
signs. A completion immediately gives \eqref{eq:factor}. Conversely,
such a factor has rank $r$ and $BB^{\mathsf T}=nI_r$, because its
nonzero Gram eigenvalues are all $n$. Its row space is the range of
$C$, perpendicular to the row space of $P$, so stacking $P$ and $B$
completes the matrix.
\end{proof}

\subsection{A finite-type reduction for each fixed defect}
\begin{proposition}[Anchor enumeration]\label{prop:anchor}
For a fixed $r$, deciding \eqref{eq:factor} reduces to at most
$2^{r^2}$ sign anchor matrices, followed by exact rational linear
algebra and sign tests on all columns. The set of possible anchored
column types is finite independently of $n$.
\end{proposition}
\begin{proof}
Choose $r$ linearly independent columns of any real Gram factor of
$C$; equivalently choose an invertible principal submatrix $G$ of $C$.
The Gram matrix of these anchors has diagonal $r$ and entries in the
finite set $\{-r,-r+2,\ldots,r\}$. A sign factor must have an
invertible $r\times r$ anchor matrix $S$ satisfying $S^{\mathsf T}S=G$.
Enumerate the $2^{r^2}$ possible $S$. For each column $j$, let $z_j$
be its $r$ correlations with the anchors. Its only possible factor
column is
\[
 b_j=(S^{\mathsf T})^{-1}z_j.
\]
Test whether every entry of every $b_j$ is $\pm1$. Since $C$ has rank
$r$, the Schur-complement identity forces these columns to have the
whole Gram matrix $C$; a final exact Gram comparison also certifies it.
There are only $(r+1)^r$ possible correlation vectors $z_j$. Both the
anchor matrices and possible column types thus come from finite lists
depending only on $r$, though their multiplicities depend on $n$.
\end{proof}

This is an elementary fixed-parameter decision reduction, not a proof
that a factor exists. It suggests importing classification of integral
quadratic forms and spherical configurations. Gram-factorization and
lattice-reduction algorithms are already used for this problem;
see Goldberger--Strassler~\cite{LLL}. No new algorithmic priority is
claimed for the reduction.

\begin{proposition}[An elementary defect-two completion]
For even $n\ge2$, every real partial Hadamard matrix with $n-2$ rows
and $n$ columns completes.
\end{proposition}
\begin{proof}
Factor its rank-two residual as columns $v_j\in\R^2$. They have squared
length two, and pairwise inner products in $\{-2,0,2\}$ by integrality,
parity and Cauchy--Schwarz. Thus any two are parallel or orthogonal.
Rank two forces exactly two perpendicular lines. Rotate the lines to
the directions $(1,1)$ and $(1,-1)$. All columns then have sign
coordinates, giving \eqref{eq:factor}. This independently recovers a
small case of the classical completion theory.
\end{proof}

\subsection{A valid-order obstruction and a successful repair}
Classical near-completion results reach seven missing rows; see
Verheiden~\cite{Verheiden}. Beyond that range arbitrary partial matrices
need not extend, and counterexamples to broad extension assertions
are documented by Bulutoglu--Kaziska~\cite{BK}. We therefore tested
completion and row exchange rather than assuming a greedy extension
lemma.

An exact seeded experiment with $200$ greedy starts at order $20$
produced $80$ maximal matrices with $12$ rows and $120$ full matrices.
These are properties of that sampling procedure, not estimates of all
equivalence classes or all partial matrices. The first terminal example
has rows
\begin{center}\small\ttfamily
\begin{tabular}{l}
++++++++++++++++++++\\
++++++++++----------\\
+++++-----+++++-----\\
++--+---++--+-++-+-+\\
++---++--+++--++--+-\\
+++-+-+---+-----++++\\
+-++----+++--+-+--++\\
+--++++-----++-+-++-\\
++-+-+---+-+-+--++-+\\
+--+++-+--+---+++--+\\
+-++--+-+--+--++++--\\
++--+--++--+-+-++-+-
\end{tabular}
\end{center}
The JSON file, rather than the display, is the authoritative certificate.
All pairwise row products are zero. No sign row is orthogonal to all
twelve: independent verification exhausts the $\binom{19}{10}=92378$
balanced vectors with first entry positive; every other vector is
excluded by the all-one row or global negation.

The rank-eight residual satisfies all PSD, integral, projection and
parity conditions above, but has no sign Gram factor. Thus those
conditions are insufficient even at order $20$, where Hadamard matrices
exist. It is a counterexample to a completion shortcut, not to the
Hadamard conjecture.

Deleting any one of the twelve rows leaves only that same row, up to
negation, as an orthogonal sign replacement; the supplied search checks
all $2^{19}$ normalized sign vectors for each deletion. Therefore this
example is frozen under one-row replacement. Deleting rows $4$ and $5$
(indices starting at one), and inserting three rows instead, gives
thirteen orthogonal rows and a full exact completion. The replacement
rows are recorded as negative-entry bitmasks $429230,635086,713014$,
with coordinate zero the least significant bit. A full order-$20$
completion is stored and its Gram matrix is independently checked.

\begin{question}[Exchange augmentation]\label{q:exchange}
Which integral equal-norm tight-frame classes admit bounded-size
exchanges that augment an orthogonal sign subframe? In the Hadamard
case, can one identify an invariant that permits repeated augmentation
without entering an unrepairable terminal configuration?
\end{question}
This asks for a broader structural exchange theorem, inspired by
matroid augmentation, lattice theory and discrete Morse theory. The
orthogonality system is not a matroid: the example above has a maximal
set of size twelve as well as sets of size twenty. A universal theorem
that every partial matrix at an admissible order can delete at most a
fixed number $c$ of rows and insert one more row than it deleted would
imply the full conjecture by iteration. The data already require
$c\ge2$; no universal finite $c$ is established here. Classifying
fixed-defect anchored configurations and their exchange costs is a
concrete intermediate problem before making that much stronger claim.

\begin{proposition}[Tensor amplification of a dead end]
If a partial sign matrix $P$ has no orthogonal sign extension, and $H$
is any square Hadamard matrix, then $P\otimes H$ also has no such
extension.
\end{proposition}
\begin{proof}
Write a putative extension as an $n\times s$ sign matrix $X$, where
$s$ is the order of $H$. Orthogonality gives $PXH^{\mathsf T}=0$ up to
the choice of vectorization convention. Since $H$ is invertible,
$PX=0$. Every column of $X$ would be a sign extension of $P$, a
contradiction.
\end{proof}
Thus arbitrary-extension failures are not isolated small-order accidents.
This does not prove any lower bound on the number of rows needed to
repair the amplified examples.

\section{Geometric optimization and its exactness threshold}
The orthogonal group gives another cross-disciplinary formulation:
\[
 F(U)=\sum_{i,j}|U_{ij}|\le n\sqrt n\qquad(U\in O(n)),
\]
with equality exactly at scaled Hadamard matrices. A naive plan is to
prove a benign optimization landscape and follow gradient ascent.
The following known basic family, independently certified here,
disproves that plan~\cite{Almost}.

\begin{proposition}[Strict local traps at admissible orders]
For every $n>4$, $U=2J_n/n-I_n$ is a strict local maximum of $F$ on
$O(n)$, but is not a scaled Hadamard matrix. Its value is $3n-4$.
\end{proposition}
\begin{proof}
$U^2=I$. Its sign matrix is $S=J_n-2I_n$, and
\[
 S^{\mathsf T}U=2I_n+(1-4/n)J_n
\]
is positive definite: its eigenvalues are $n-2$ and $2$. In the open
sign chamber containing $U$, $F(V)=\operatorname{tr}(S^{\mathsf T}V)$.
Writing $S=UM$ with $M>0$, the orthogonal Procrustes inequality gives
$\operatorname{tr}(M U^{\mathsf T}V)\le\operatorname{tr}M$, with
equality only for $V=U$. To see uniqueness directly, diagonalize $M$:
every diagonal entry of an orthogonal matrix is at most one, and all
weights are strictly positive. Finally direct summation gives $3n-4$;
the diagonal and off-diagonal magnitudes differ when $n>4$.
\end{proof}
Tensor products preserve these strict-local-maximum certificates:
the sign matrices and positive polar factors tensor. Thus tensoring
the order-eight trap produces traps at arbitrarily large admissible
orders. This rejects a specific landscape claim, not every possible
continuous objective or nonlocal optimization method.

\begin{proposition}[How accurate an approximation would suffice]
For any non-Hadamard sign matrix $A$ of even order $n>2$,
\[
 \operatorname{cond}_2(A)\ \ge\sqrt{\frac{n+2}{n-2}},
\]
where singular matrices have infinite condition number.
\end{proposition}
\begin{proof}
A nonzero off-diagonal Gram entry has magnitude at
least two. Its $2\times2$ principal submatrix has eigenvalues
$n\pm|g|$, so interlacing gives the bound for nonsingular $A$.
The singular case is immediate.
\end{proof}

For the intended even orders, the separation is $1+2/n+O(n^{-2})$.
Thus a sign construction with condition number $1+o(1/n)$ would force
eventual exact Hadamard existence. A bound merely tending to one does
not. An eventual result would still need a rigorous finite treatment
below its threshold to prove the full conjecture.

\section{Why a local border cannot implement order induction}
Absorption and trade arguments from design theory are attractive, but
their local repair must respect a global rank constraint.
\begin{proposition}[Preserved-rectangle obstruction]\label{prop:border}
Let $H_n$ and $K_{n+s}$ be Hadamard matrices, $s>0$. Suppose a
$k\times(n-c)$ rectangle of $H_n$ appears unchanged in $K_{n+s}$.
Then $k\le c+s$. In particular preserving an $(n-r)\times(n-c)$
rectangle requires $r+c\ge n-s$.
\end{proposition}
\begin{proof}
After reordering, write the selected old rows as $[A\ B]$, with
$B$ having $c$ columns, and the corresponding new rows as $[A\ D]$,
with $D$ having $c+s$ columns. Orthogonality says
\[
 AA^{\mathsf T}+BB^{\mathsf T}=nI_k,\qquad
 AA^{\mathsf T}+DD^{\mathsf T}=(n+s)I_k.
\]
Hence $DD^{\mathsf T}=sI_k+BB^{\mathsf T}$ is positive definite,
and $k=\operatorname{rank}D\le c+s$.
\end{proof}
For example, adjoining four rows and columns while preserving almost
all old rows and columns cannot work. An additive-order construction
would need changes spread through a large part of the old matrix.
This is stronger information than a failed numerical border search.

The existence theorems proved by iterative absorption concern a fixed
block hypergraph with sufficiently large host size~\cite{Absorption}.
The equivalent Hadamard design has parameters
$2$-$(4m-1,2m-1,m-1)$: its block size grows with its point set.
Those fixed-block theorems therefore cannot be applied at these
parameters without new uniform estimates and exact-support control.
An absorption theorem for growing blocks at this tight boundary would
be a substantial theorem in its own right; it is not available here.

\section{Priorities after this first research pass}
The first priority is Question~\ref{q:norm}: analyze integral
congruence-constrained Hermitian representation over the cyclotomic
components of $R_m$, including compatibility at primes dividing $m$.
An independently interesting larger theorem would characterize which
self-adjoint positive elements of an integral group ring are sums of
four Hermitian norms in prescribed residue classes. Our scalar target
would then be one specialization. Positivity and local congruence
solvability must not be assumed sufficient without a proof; integral
representation obstructions at small energy are exactly the issue.
First seek obstructions and sufficient criteria above minimum energy,
not a claimed local-to-global theorem. The boundary requires simultaneous
control at all characters. Full quaternionic perfection has already
been rejected by published counterexamples.

The second priority is the integer moment problem: use
Propositions~\ref{prop:signed}--\ref{prop:moves} to study conversion of
signed designs into nonnegative ones. Identify an obstruction invariant
or an elimination theorem beyond generic normality. The nonnormality
example and the rank bound must be built into any proposed theorem.
The central boundary case alone is equivalent to the original conjecture.

The third priority is a classification of residual frames at fixed
defect, with row-exchange costs. Proposition~\ref{prop:anchor} makes
this finite at the level of column types. The certified order-twenty
example is a useful adversarial test: direct extension and one-row
replacement fail, while two-for-three exchange succeeds. This route
still needs a way to reach a near-complete frame at every order; a
terminal repair theorem alone cannot do that.

We have proved reductions and negative tests, not a promising complete
proof. The note's value is a set of precise, falsifiable intermediate
questions and independently checked certificates. A list of connected
subjects, an approximate solution, or a finite search success is not
being presented as evidence that the full conjecture is proved.

\section{Reproduction and scope}
All files are under \path{research/full_conjecture/}. The scripts use
Python's standard library only. No external research code, supplement
or Lean proof was imported. Run \path{experiments/investigate.py},
\path{experiments/exchange.py}, \path{experiments/moments.py},
\path{experiments/norm_lifting.py}, then \path{experiments/verify.py}.
The independent verifier checks the dead end, its augmentation and
completion, and the sequence norm certificates using exact arithmetic.
The kernel-basis and minimum-energy assertions have the paper proofs
above; no Lean formalization or novelty claim is made.

\begin{thebibliography}{10}
\small\raggedright
\bibitem{Survey} R. Egan, I. Kotsireas, S. London and P. \'O Cath\'ain,
\emph{Recent progress in the construction of Hadamard matrices}, 2026.
\url{https://arxiv.org/abs/2610.27175}.
\bibitem{Quaternion} A. Bennett, C. Bright, P. Colinot and A. Nayak,
\emph{Quaternionic Perfect Sequences and Hadamard Matrices}, 2026,
Theorem 9 and introduction. \url{https://arxiv.org/abs/2601.22337}.
\bibitem{KLP} G. Kuperberg, S. Lovett and R. Peled,
\emph{Probabilistic existence of regular combinatorial structures},
version 3, 2017. \url{https://arxiv.org/abs/1302.4295}.
\bibitem{Davis} D. Davis, \emph{Counting partial Hadamard matrices in the
cubic regime}, 2026, Theorem 1.1. \url{https://arxiv.org/abs/2603.30013}.
\bibitem{Ohsugi} H. Ohsugi, \emph{Normality of cut polytopes of graphs
is a minor closed property}, Discrete Math. 310 (2010), 1160--1166.
\url{https://arxiv.org/abs/0906.5303}.
\bibitem{LLL} A. Goldberger and Y. Strassler, \emph{A practical algorithm
for completing half-Hadamard matrices using LLL}, J. Algebraic Combin.
55 (2022), 217--244. Related authors' presentation:
\url{https://www.netanya.ac.il/gur/data/Hadamard_Research/research9.pdf}.
\bibitem{Verheiden} E. P. Verheiden, \emph{Arithmetical Properties of
Combinatorial Matrices}, Ph.D. thesis, Caltech, 1978.
\url{https://thesis.caltech.edu/2899/}.
\bibitem{BK} D. A. Bulutoglu and D. M. Kaziska, \emph{A counterexample
to Beder's conjectures about Hadamard matrices}, J. Statist. Plann.
Inference 139 (2009), 3381--3383.
\url{https://doi.org/10.1016/j.jspi.2009.02.003}.
\bibitem{Almost} T. Banica, I. Nechita and K. \.Zyczkowski,
\emph{Almost Hadamard matrices: general theory and examples},
Open Syst. Inf. Dyn. 19 (2012), 1250024.
\url{https://arxiv.org/abs/1202.2025}.
\bibitem{Absorption} S. Glock, D. K\"uhn, A. Lo and D. Osthus,
\emph{The existence of designs via iterative absorption: hypergraph
$F$-designs for arbitrary $F$}, version 3.
\url{https://arxiv.org/abs/1611.06827}.
\bibitem{Filters} S.-M. Phoong and Y.-P. Lin,
\emph{Lapped Hadamard transforms and filter banks},
Proc. ICASSP 2003, vol. 6, 509--512.
\url{https://doi.org/10.1109/ICASSP.2003.1201730}.
\bibitem{Brock} B. W. Brock, \emph{Hermitian congruence and the existence
and completion of generalized Hadamard matrices}, J. Combin. Theory
Ser. A 49 (1988), 233--261.
\url{https://doi.org/10.1016/0097-3165(88)90054-4}.
\end{thebibliography}
\end{document}
